\chapter{تنظيم نظام التشغيل}
\label{CH:FIRST}

تتمثل إحدى المتطلبات الأساسية في تصميم أنظمة التشغيل في تحقيق \textbf{العزل القوي} (strong isolation).
يقتضي العزل منع العمليات من إتلاف أو التجسس على ذاكرة وحالات العمليات الأخرى، وكذلك حماية نظام التشغيل نفسه من البرامج الخبيثة أو المعطوبة التي تعمل في فضاء المستخدم.

يقدم هذا الفصل كيفية تنظيم نظام التشغيل لتحقيق هذا العزل، والتعامل مع وضعيات تنفيذ المعالج، وطريقة إقلاع النواة وإنشاء أول عملية في xv6.

\section{تجريد الموارد الفيزيائية}

أحد الطرق البدائية لتحقيق التفاعل مع الموارد هي السماح للبرامج بالوصول المباشر لعتاد الحاسوب (الذاكرة، القرص، وحدة المعالجة). لكن هذا النهج يفشل في تقديم العزل؛ إذ يمكن لبرنامج معطوب أن يكتب فوق ذاكرة برنامج آخر أو يسيطر على العتاد.

لذلك، ينشئ نظام التشغيل \textbf{تجريدات محكمة} (abstractions) فوق العتاد المادي:
\begin{itemize}
\item بدلاً من السماح للبرامج بالوصول المباشر للذاكرة الفيزيائية، يوفر نظام التشغيل لكل عملية \indextext{فضاء ذاكرة ظاهري} (virtual address space) خاصًا بها.
\item بدلاً من الوصول المباشر للقرص الصلب، يقدم نظام التشغيل تجريد \indextext{نظام الملفات} (file system).
\item بدلاً من التحكم المباشر بأنوية المعالج، يوفر نظام التشغيل تجريد \indextext{العملية} (process) والجدولة الزمنية.
\end{itemize}

\section{وضع المستخدم، وضع المشرف، واستدعاءات النظام}

يتطلب تحقيق العزل القوي دعمًا صريحًا من عتاد وحدة المعالجة المركزية (CPU).
توفر معمارية \textbf{RISC-V} ثلاثة مستويات من الصلاحيات والأنماط التي ينفذ فيها المعالج التعليمات:
\begin{enumerate}
\item \textbf{وضع الآلة (Machine Mode - M-mode)}: أعلى مستوى صلاحية، يمتلك وصولاً كاملاً وغير مشروط للعتاد، ويُستخدم عادة لتهيئة العتاد عند الإقلاع الأول (مثل البرمجيات الثابتة Firmware / BIOS).
\item \textbf{وضع المشرف (Supervisor Mode - S-mode)}: المستوى المخصص لنواة نظام التشغيل. يسمح بتنفيذ التعليمات الممتازة (privileged instructions) مثل تعديل سجلات إدارة الذاكرة وتفعيل المقاطعات.
\item \textbf{وضع المستخدم (User Mode - U-mode)}: المستوى المخصص لبرامج المستخدم العادية. يُمنع المعالج في هذا الوضع من تنفيذ التعليمات الممتازة أو الوصول المباشر للسجلات العتادية الحساسة.
\end{enumerate}

عندما يحاول برنامج يعمل في وضع المستخدم تنفيذ تعليمة ممتازة (مثل تعديل سجل التحكم بالتصفح \lstinline{satp})، يرفض العتاد التعليمة ويولّد \indextext{استثناء} (exception) يُحاط به في النواة.

للإنتقال من وضع المستخدم إلى وضع المشرف بأمان، يوفر عتاد RISC-V تعليمة خاصة تُدعى \lstinline{ecall} (Environment Call).
عند تنفيذ \lstinline{ecall}:
\begin{enumerate}
\item يحفظ المعالج حالة البرنامج ويعيد رفع مستوى الصلاحية إلى وضع المشرف (S-mode).
\item يقفز المعالج إلى عنوان معالج المصايد (trap handler) المسجل مسبقًا في السجل الممتاز \lstinline{stvec}.
\item تنفذ النواة الخدمة المطلوبة وتتحقق من صحة المدخلات.
\item تعود النواة إلى وضع المستخدم باستخدام تعليمة \lstinline{sret} (Supervisor Return).
\end{enumerate}

\section{تنظيم النواة}

تتوزع أنظمة التشغيل في هيكلة النواة إلى نمطين رئيسين:
\begin{itemize}
\item \textbf{النواة المتجانسة (Monolithic Kernel)}:
تنفذ جميع مكونات نظام التشغيل (إدارة الذاكرة، الجدولة، برامج تشغيل الأجهزة، نظام الملفات) داخل فضاء النواة وبصلاحيات كاملة (S-mode). يوفر هذا التصميم أداءً عاليًا وسرعة في استدعاءات الخدمات. يعتمد نظام \textbf{xv6}، مثل Linux و FreeBSD، هيكلة النواة المتجانسة.
\item \textbf{النواة الدقيقة (Microkernel)}:
تنفذ النواة الحد الأدنى من الخدمات (الجدولة، الاتصال بين العمليات IPC، إدارة الذاكرة المباشرة)، بينما تنفذ بقية الخدمات مثل نظام الملفات وبرامج تشغيل الأجهزة كعمليات عادية في فضاء المستخدم. يزيد هذا التصميم من أمان النظام واستقراره لكنه يتطلب تكلفة إضافية في تبديل السياق والاتصالات.
\end{itemize}

\section{الكود البرمجي: هيكلة ونواة xv6}

تتواجد الشفرة المصدرية لنواة xv6 داخل دليل \lstinline{kernel/}.
وتُقسم إلى وحدات برمجية واضحة:
\begin{itemize}
\item \lstinline{kernel/main.c}: نقطة بداية تشغيل النواة عند الإقلاع.
\item \lstinline{kernel/proc.c}: إدارة العمليات والجدولة وتغيير السياق.
\item \lstinline{kernel/vm.c}: إدارة الذاكرة الافتراضية وجداول الصفحات.
\item \lstinline{kernel/trap.c}: معالجة المصايد، الاستثناءات، واستدعاءات النظام.
\item \lstinline{kernel/spinlock.c}: أقفال التناوب للتزامن وحماية هياكل البيانات.
\item \lstinline{kernel/bio.c} و \lstinline{kernel/fs.c}: طبقات المخزن المؤقت ونظام الملفات.
\end{itemize}

\section{نظرة عامة على العمليات}

تُمثَّل كل عملية في xv6 بواسطة هيكل بيانات C يُدعى \lstinline{struct proc} (المعرف في \lstinline{kernel/proc.h}).
يحتفظ هذا الهيكل بجميع معلومات حالة العملية:
\begin{itemize}
\item حالة العملية (\lstinline{enum procstate}): هل هي \lstinline{UNUSED}، \lstinline{EMBRYO}، \lstinline{SLEEPING}، \lstinline{RUNNABLE}، \lstinline{RUNNING}، أو \lstinline{ZOMBIE}.
\item معرّف العملية (\lstinline{pid}).
\item جدول الصفحات الخاص بالعملية (\lstinline{pagetable_t pagetable}).
\item المكدس النواتي الخاص بالعملية (\lstinline{uint64 kstack}).
\item إطار المصيدة (\lstinline{struct trapframe *trapframe}) للحفظ والاسترجاع لسجلات المستخدم.
\item سياق العملية (\lstinline{struct context context}) لحفظ سجلات النواة أثناء تبديل المجدول.
\end{itemize}

\section{الكود البرمجي: إقلاع xv6، العملية الأولى واستدعاء النظام}

عند تشغيل الحاسوب (أو محاكي QEMU):
\begin{enumerate}
\item يبدأ المعالج في تنفيذ كود الـ ROM الإقلاعي في وضع الآلة (M-mode)، والذي يحمل النواة إلى الذاكرة عند العنوان الفيزيائي \lstinline{0x80000000}.
\item ينفذ المعالج كود التجميع في \lstinline{kernel/entry.S}، لتهيئ مكدس النواة الابتدائي لكل نوات معالجة (hart).
\item تنتقل السيطرة إلى الدالة \lstinline{main()} في \lstinline{kernel/main.c} في وضع المشرف (S-mode).
\item تقوم النواة بتهيئة الأقفال، ووحدة إدارة الذاكرة، ومجمع المقاطعات PLIC، والذاكرة المخبئية.
\item تدعو \lstinline{main()} الدالة \lstinline{userinit()} لإنشاء أول عملية مستخدم.
\item تنفذ العملية الأولى كوداً مجمعاً صغيراً بداخل \lstinline{usercnt.S} يتضمن استدعاء النظام \lstinline{exec("/init")} ليعمل برنامج البدء الرسمي الشارح لمغلف الأوامر \lstinline{sh}.
\end{enumerate}

\section{نموذج الأمان}

يفترض نموذج الأمان في xv6 أن النواة موثوقة تمامًا ومحصنة ضد الثغرات، بينما برامج المستخدم غير موثوقة ومفصولة تمامًا. تقوم النواة بالتحقق من جميع العناوين الممررة من برامج المستخدم للتأكد من وقوعها ضمن النطاق المسموح به للعملية قبل إجراء القراءة أو الكتابة.

\section{العالم الحقيقي}

تستخدم أنظمة التشغيل الإنتاجية مثل Linux آليات أمان متطورة مثل مساحات الأسماء (Namespaces) والحاويات (Containers) وASLR (عشوائية تخطيط فضاء العناوين) لتعزيز العزل. ومع ذلك، تشترك جميع هذه الأنظمة مع xv6 في الأعمدة الثنائية ذاتها: العزل العتادي بين وضع المستخدم ووضع المشرف، واستخدام استدعاءات النظام المحكومة عبر التعليمة \lstinline{ecall}.

\section{تمارين}

1. تتبع خطوات تنفيذ تعليمة \lstinline{ecall} بدءًا من وضع المستخدم وصولاً إلى معالج \lstinline{syscall()} في \lstinline{kernel/syscall.c}.
2. أضف استدعاء نظام جديد باسم \lstinline{trace(mask)} يقوم بطباعة أسماء استدعاءات النظام المنفذة لاحقاً.
